\documentclass[11pt]{article}

\usepackage[margin=1in]{geometry}
\usepackage{amsmath}
\usepackage{amssymb}

\usepackage{graphicx}
\usepackage{bm}
\usepackage{booktabs}
\usepackage{array}
\usepackage{amsthm}
\usepackage{enumerate}
\usepackage[numbers,sort&compress]{natbib}
\usepackage{hyperref}
\usepackage{lineno}

\hypersetup{
  colorlinks=true,
  linkcolor=blue,
  citecolor=blue,
  urlcolor=blue
}

\newtheorem{theorem}{Theorem}
\newtheorem{proposition}{Proposition}
\newtheorem{corollary}{Corollary}
\newtheorem{lemma}{Lemma}

\begin{document}

\title{The Universal Mass Sum Rule: A Theta-Independent Prediction for All CDFT Torus Knot Families}

\author{Steve Bico Mujjabi, MD\thanks{msbico@gmail.com}\\
VuraLabs, Kampala, Uganda\\
Independent researcher\\
ORCID: \href{https://orcid.org/0009-0001-0556-5516}{0009-0001-0556-5516}\\
\url{https://github.com/bampita-bico/CDFD-Part-I-Release}}

\date{July 2026}

\maketitle

\begin{abstract}
Paper VI argued within the model that the Brannen amplitude coefficient is
universal across all Z-n torus-knot families and that the Z-three family is
uniquely exactly solvable. For family numbers of five or greater, the
individual mode masses depend on an orientation angle that cannot be fixed by
pure algebra alone.

The present paper identifies what can still be derived without knowing that
orientation. Its central result is the universal mass sum rule: for every
T(2,n) family, the total mass of the family equals twice the family number
times the family energy scale. The proof follows directly from the Z-n Fourier
identity and does not depend on the orientation angle. Combined with
Faddeev-Niemi scaling, this means the total family mass grows as the family
number to the seven-fourths power.

For the trefoil family, the rule reproduces the observed charged-lepton total
within the tau-mass uncertainty reported in Paper V. For the cinquefoil family,
it gives a total mass of about 4604 MeV. The paper also proves the Z-five mode
split: at the symmetric point, the five amplitudes divide into three positive
and two negative modes. A unique Z-five self-consistency angle near
3.70 degrees is derived and reported as a hypothesis pending a first-principles
physical argument.

Across the series, the public CDFD Runtime expresses this equilibrium vacuum
limit as a reusable flow-constraint state grammar. The calculations below use
that equilibrium limit only; adaptive departures are left for later dynamical
work rather than fitted in this manuscript. The paper-local supplementary
scripts remain the audited reproduction layer for the figures and tables in
this Part I archive.
\end{abstract}

\noindent\textbf{Keywords:} Constraint-Driven Field Theory, Constraint-Driven Flux Dynamics, adaptive operating ratio, vacuum structure, Mujjabi tests, reproducible physics, mass sum rule, particle spectra

\vspace{0.75em}
\linenumbers

\section*{Project Context}
This manuscript constitutes Part I (Paper 07) of the multi-part series
\textit{Constraint-Driven Flux Dynamics (CDFD)}. The underlying mathematical
architecture builds directly upon the concepts established in Paper 06 of this
series~\cite{MujjabiVI2026}.

\section*{Symbol Conventions}

\noindent\textbf{CDFD physical-vacuum symbol convention.}
\begin{center}
{\small
\begin{tabular}{p{0.16\linewidth}p{0.33\linewidth}p{0.41\linewidth}}
\hline
Symbol & Part I meaning & Physical reading \\
\hline
$J$ or $\Phi$ & Driving flux or throughput & Transport current, field throughput, circulation drive, or generic flux; a subscript fixes the exact channel \\
$C$ & Active constraint or capacity burden & Vacuum transport capacity, geometric confinement, stiffness, boundary resistance, or load-bearing limit \\
$S$ & Surface responsiveness & Local ability of the vacuum or modeled surface to reconfigure under drive \\
$M_s$ & Structural memory & Retained geometry, phase history, stiffness history, or accumulated constraint state \\
$\Psi_s$ & Adaptive operating ratio $(J/C)\cdot S\cdot M_s$ or $(\Phi/C)\cdot S\cdot M_s$ & Local bookkeeping gauge for constrained operation; it is not a new universal constant \\
$R,a,\chi$ & Ring radius, core radius, and aspect ratio $R/a$ & Geometry used to connect vortex structure to dimensionless electromagnetic scales \\
$\rho_0,c,\kappa$ & Vacuum density scale, propagation speed, and transport stiffness & Medium parameters used in stability, pressure, and self-consistency calculations \\
\hline
\end{tabular}
}
\end{center}

\noindent\textbf{Greek-symbol convention.}
\begin{center}
{\small
\begin{tabular}{p{0.16\linewidth}p{0.33\linewidth}p{0.41\linewidth}}
\hline
Symbol & Local role & Interpretation rule \\
\hline
$\alpha$ & Fine-structure constant unless a local equation states otherwise & In the Part I physics papers, un-subscripted $\alpha$ means the electromagnetic fine-structure constant \\
$\beta,\gamma$ & Local response, stiffness, coupling, or correction coefficients & Equation-local coefficients; subscripts bind them to the named mode, channel, or approximation \\
$\theta,\phi,\Phi_c$ & Phase, orientation, or chirality angles & Used for torus-knot orientation, vacuum phase, or chiral phase geometry as defined in the nearby equation \\
$\delta,\epsilon$ & Perturbation, residual, tolerance, or small-offset scales & Local stability margins, numerical residuals, or experimental tolerances \\
$\lambda,\mu,\nu$ & Rate, mass, mode, or frequency labels & Meaning is fixed by the equation, subscript, and physics context where the symbol appears \\
$\rho,\sigma,\tau$ & Density, spread/noise, or time-scale terms & Local density, variance, stress, relaxation, or transport-memory quantities \\
$\Omega,\Lambda$ & State/domain space or integrated measure where defined & Reserved for explicitly defined state spaces, spectra, or synthesis-level quantities \\
\hline
\end{tabular}
}
\end{center}


\section{Introduction}

The CDFT series has established a complete derivation chain for the lepton
sector (Papers~I--V~\cite{MujjabiI2026,MujjabiII2026,MujjabiIII2026,MujjabiIV2026,MujjabiV2026})
and identified the universal coefficient $c_n = \sqrt{2}$ and the
$\mathbb{Z}_3$ uniqueness theorem (Paper~VI~\cite{MujjabiVI2026}).

A key finding of Paper~VI is that for $n \geq 5$ the Brannen amplitudes
$A_k = 1 + \sqrt{2}\cos(\theta + 2\pi k/n)$ are negative for at least one
mode at every orientation $\theta$.
This structural obstruction means that the Koide-like invariant $Q_n = 2/n$
(which requires all $A_k \geq 0$) does not hold as a constant, and the
purely algebraic self-consistency argument that determined $\theta_3 = 2/9$~rad
in Paper~V cannot be directly applied.

The question for the present paper is: \emph{what predictions can be made
for $n \geq 5$ families without knowing $\theta_n$?}

The answer is the universal mass sum rule: the sum of all mode masses is
$2nM_n$, exactly, for any $\theta_n$.
This follows from a two-line algebraic identity and requires no knowledge
of $\theta_n$.
For $n = 3$ it recovers the known lepton sum $m_e + m_\mu + m_\tau$ to
within the Paper~V tau mass error.
For $n = 5$ it predicts the total cinquefoil mass sum as $4603.85$~MeV.

\subsection{Mass as constrained transport cost}

In the series architecture, this paper is where mass becomes a budget rather
than a label. The Mujjabi Geometric Charge Principle asks how a boundary
discontinuity is measured; the corresponding mass reading is that inertia is the
transport cost of maintaining a stable constrained vortex state. The exact sum
rule
\[
  \sum_k m_k=2nM_n
\]
then says that each family has a conserved action/transport budget even when
the individual mode allocation depends on $\theta_n$.

This is also where the open Mujjabi Action Principle belongs. A complete theory
would derive the action quantum from a critical flux threshold and vortex
period, $h=\Phi_c\tau_v$, rather than using $\hbar$ as input. The present paper
does not complete that derivation; it states the mass-budget structure that such
a derivation must explain.


\section{The Universal Mass Sum Rule}
\label{sec:sum_rule}

\subsection{Statement and proof}

\begin{theorem}[Universal mass sum rule]
  \label{thm:sum_rule}
  For any $T(2,n)$ family with masses $m_k = M_n A_k^2$,
  $A_k = 1 + \sqrt{2}\cos(\theta + 2\pi k/n)$, $k = 0,\ldots,n-1$:
  \begin{equation}
    \sum_{k=0}^{n-1} m_k \;=\; 2\,n\,M_n\,,
    \label{eq:sum_rule}
  \end{equation}
  for any $\theta$ and any $M_n > 0$.
\end{theorem}

\begin{proof}
  From the $\mathbb{Z}_n$ Fourier identity (Paper~VI, Eq.~(2)):
  $\sum_{k=0}^{n-1} \cos^2(\theta + 2\pi k/n) = n/2$ for any $\theta$.
  Therefore:
  \[
    \sum_k A_k^2
    = \sum_k \bigl[1 + \sqrt{2}\cos(\theta + 2\pi k/n)\bigr]^2
    = n + 2\sqrt{2}\cdot 0 + 2\cdot\frac{n}{2} = 2n\,.
  \]
  Hence $\sum_k m_k = M_n \sum_k A_k^2 = 2nM_n$.
\end{proof}

\begin{corollary}[Mean mode mass]
  The mean mass per mode is $\bar{m} = 2M_n$, independent of $\theta_n$.
\end{corollary}

\subsection{Why this is non-trivial}

At first glance, Eq.~\eqref{eq:sum_rule} might seem like a restatement of
known identities.
The non-trivial content is twofold:
\begin{enumerate}[(a)]
  \item The sum is \emph{theta-independent}.
    Individual masses $m_k = M_n A_k^2$ depend strongly on $\theta_n$:
    they range from $0$ (when $A_k = 0$) to $M_n(1+\sqrt{2})^2 \approx 5.83 M_n$.
    Yet their sum is always exactly $2nM_n$.
    This means the total mass of a family is a firm prediction even when
    the mass spectrum (individual mode masses) is not.

  \item Combined with the Faddeev--Niemi scaling $M_n = M_3(n/3)^{3/4}$
    (Paper~IV~\cite{MujjabiIV2026}), the sum becomes a prediction from
    first principles:
    \begin{equation}
      \sum_{k=0}^{n-1} m_k(n)
      \;=\; 2n M_3 \left(\frac{n}{3}\right)^{3/4}
      \;=\; 2M_3 \cdot 3^{-3/4} \cdot n^{7/4}\,.
      \label{eq:sum_n74}
    \end{equation}
    This requires only $M_3$ (derived from $m_e$ in Paper~V) and the
    Faddeev--Niemi exponent $3/4$ (from the field theory of knotted
    solitons~\cite{FaddeevNiemi1997}).
    No other free parameters enter.
\end{enumerate}

\subsection{Verification for $n = 3$}

For the lepton family ($n = 3$, $M_3 = 313.86$~MeV):
\begin{align}
  \sum_k m_k^{(3)} &= 6 M_3 = 1883.16~\text{MeV}\,.
  \label{eq:sum_3}
\end{align}
The measured value is $m_e + m_\mu + m_\tau = 1883.03$~MeV (PDG~2022~\cite{PDG2022}).
The difference is $0.13$~MeV, which equals the tau mass prediction error
from Paper~V ($+1.04\sigma$).
The sum rule is verified at the accuracy of the Paper~V prediction.


\section{The $n^{7/4}$ Mass Hierarchy}
\label{sec:hierarchy}

\subsection{Derivation of the scaling law}

From Eq.~\eqref{eq:sum_n74}, the total family mass scales as:
\begin{equation}
  \Sigma_n \;\equiv\; \sum_{k=0}^{n-1} m_k(n) \;=\; 2M_3 \cdot 3^{-3/4} \cdot n^{7/4}\,.
  \label{eq:n74}
\end{equation}
The exponent $7/4 = 1.75$ arises as:
\begin{align*}
  \Sigma_n &= 2n \cdot M_3 \left(\frac{n}{3}\right)^{3/4}
  = 2M_3 \cdot 3^{-3/4} \cdot n^{1+3/4}
  = 2M_3 \cdot 3^{-3/4} \cdot n^{7/4}\,.
\end{align*}
The $n$ factor comes from the number of modes (the sum has $n$ terms).
The $n^{3/4}$ factor comes from Faddeev--Niemi energy scaling.
Their product gives $n^{7/4}$.

\subsection{Hierarchy table}

\begin{table}[h]
  \caption{CDFT torus knot mass sum hierarchy.
    All predictions from Theorem~\ref{thm:sum_rule} and Faddeev--Niemi
    scaling, zero free parameters beyond $m_e$.}
  \label{tab:hierarchy}
  \begin{tabular}{ccccc}
    \toprule
    $n$ & Knot & $M_n$ (MeV) & $\Sigma_n$ (MeV) & Status \\
    \midrule
    3  & $T(2,3)$ trefoil    & 313.9  & 1\,883.2  & Verified (PDG: 1\,883.0) \\
    5  & $T(2,5)$ cinquefoil & 460.4  & 4\,603.9  & Predicted \\
    7  & $T(2,7)$ heptafoil  & 592.5  & 8\,295.6  & Predicted \\
    9  & $T(2,9)$            & 715.4  & 12\,878.0 & Predicted \\
    11 & $T(2,11)$           & 831.6  & 18\,296.2 & Predicted \\
    13 & $T(2,13)$           & 942.7  & 24\,509.0 & Predicted \\
    \bottomrule
  \end{tabular}
\end{table}

\begin{figure}[ht]
  \centering
  \includegraphics[width=0.8\textwidth]{figures/fig7_0_panel.pdf}
  \caption{The mass sum hierarchy $\Sigma_n$. The total mass of the $n$-th family scales exactly as $n^{7/4}$, driven by the combination of the topological mode count $n$ and the Faddeev-Niemi field energy $n^{3/4}$.}
  \label{fig:mass_sum_hierarchy}
\end{figure}

The $n^{7/4}$ scaling is the universal mass growth law of the CDFT
torus knot hierarchy.
It is a direct consequence of two independently derived results:
the universal Brannen coefficient $c_n = \sqrt{2}$ (Paper~VI) and
the Faddeev--Niemi energy bound (Paper~IV).


\section{The $\mathbb{Z}_5$ Mode Split}
\label{sec:Z5_split}

\subsection{The $3 + 2$ pattern at the symmetric point}

\begin{theorem}[$\mathbb{Z}_5$ mode split]
  \label{thm:Z5_split}
  At the $\mathbb{Z}_5$ symmetric point $\theta = 0$, the five cinquefoil
  amplitudes divide as:
  \begin{itemize}
    \item 3 \emph{positive-phase} modes: $A_k > 0$ (aligned with vacuum background),
    \item 2 \emph{negative-phase} modes: $A_k < 0$ (anti-aligned).
  \end{itemize}
\end{theorem}

\begin{proof}
  At $\theta = 0$: $A_k = 1 + \sqrt{2}\cos(2\pi k/5)$.
  The five cosines are
  $\cos(0^\circ) = 1$,
  $\cos(72^\circ) = \cos(288^\circ) \approx 0.309$,
  $\cos(144^\circ) = \cos(216^\circ) \approx -0.809$.
  The corresponding amplitudes are:
  \begin{align*}
    A_0 &= 1 + \sqrt{2} \approx 2.414 > 0 \\
    A_1 = A_4 &= 1 + \sqrt{2}\cos72^\circ \approx 1.437 > 0 \\
    A_2 = A_3 &= 1 + \sqrt{2}\cos144^\circ \approx -0.144 < 0
  \end{align*}
  so 3 positive and 2 negative.
\end{proof}

\begin{figure}[ht]
  \centering
  \includegraphics[width=0.8\textwidth]{figures/fig7_1_panel.pdf}
  \caption{The $\mathbb{Z}_5$ mode split at $\theta = 0$. Three modes align positively with the vacuum background, while two modes oscillate in anti-phase (negative amplitude).}
  \label{fig:mode_split}
\end{figure}

\subsection{Physical interpretation of the mode split}

The negative amplitudes at $\theta = 0$ have magnitude $\sqrt{2} - 1 \approx 0.414$
below the vacuum background.
These two modes oscillate in anti-phase with the vacuum field.
All five modes have positive mass: $m_k = M_5 A_k^2 \geq 0$ always,
because mass involves the square.

For general $\theta$, the count of negative amplitudes is:
\begin{itemize}
  \item 2 negative modes for $\theta \in [0, \pi/20) \cup (7\pi/20, 2\pi/5)$,
  \item 1 negative mode for $\theta \in (\pi/20, 7\pi/20)$.
\end{itemize}
where $\pi/20 = 9^\circ$ and $7\pi/20 = 63^\circ$ are exact boundaries
(zeros of the amplitude condition $1 + \sqrt{2}\cos(\theta + 2\pi k/5) = 0$).
The 2-negative regime occupies 25\% of the period, the 1-negative regime 75\%.

The $3 + 2$ pattern at the symmetric point suggests a natural structural
division of the cinquefoil modes: three modes analogous to the three lepton
generations, and two additional modes without a $\mathbb{Z}_3$ analogue.


\section{The $\mathbb{Z}_5$ Self-Consistency Equation}
\label{sec:Z5_sc}

\subsection{The transcendental equation}

For $\mathbb{Z}_3$, the self-consistency condition $3\theta_3 = Q_3 = 2/3$
was algebraic: $Q_3$ is constant, so the solution $\theta_3 = 2/9$~rad
follows immediately (Paper~V).

For $\mathbb{Z}_5$, the modified Koide ratio (using $|A_k|$ for negative
amplitudes):
\begin{equation}
  \tilde{Q}_5(\theta) \;=\; \frac{\sum_k A_k^2}{\bigl(\sum_k |A_k|\bigr)^2}
  \;=\; \frac{10}{\bigl(\sum_k |A_k|\bigr)^2}
  \label{eq:Q5_tilde}
\end{equation}
is \emph{not} constant: it ranges in $[0.294, 0.328]$ as $\theta$ varies.
The numerator $\sum A_k^2 = 10$ is fixed (Theorem~\ref{thm:sum_rule});
only the denominator $(\sum |A_k|)^2$ varies.

The $\mathbb{Z}_5$ analogue of the chirality equilibrium (Paper~III) is:
\begin{equation}
  5\theta_5 + \phi_{c,5} \;=\; \pi\,.
  \label{eq:equil_5}
\end{equation}
The $\mathbb{Z}_5$ self-consistency condition --- that the chirality breaking
angle equals the Koide-like ratio at that angle --- then reads:
\begin{equation}
  5\theta_5 \;=\; \tilde{Q}_5(\theta_5)\,.
  \label{eq:sc_5}
\end{equation}

\begin{figure}[ht]
  \centering
  \includegraphics[width=0.8\textwidth]{figures/fig7_2_panel.pdf}
  \caption{The $\mathbb{Z}_5$ self-consistency fixed point. The transcendental equation $5\theta_5 = \tilde Q_5(\theta_5)$ has exactly one solution in the physical domain.}
  \label{fig:fixed_point}
\end{figure}

Unlike the $\mathbb{Z}_3$ case, this is a \emph{transcendental fixed-point
equation}.

\subsection{Unique solution}

\begin{proposition}[Existence and uniqueness of $\theta_5$]
  Equation~\eqref{eq:sc_5} has exactly one solution in $[0, \pi/5)$.
\end{proposition}

\begin{proof}[Proof sketch]
  At $\theta = 0$: $5\theta = 0 < \tilde{Q}_5(0) \approx 0.322$.
  At $\theta = \pi/5$: $5\theta = \pi \gg \tilde{Q}_5(\pi/5) \approx 0.294$.
  By continuity, at least one solution exists (intermediate value theorem).
  Numerically, $5\theta$ and $\tilde{Q}_5(\theta)$ cross exactly once
  (one sign change in $5\theta - \tilde{Q}_5(\theta)$ on $[0, \pi/5]$).
\end{proof}

The numerical solution to machine precision:
\begin{align}
  \theta_5 &= 0.0645346670\ldots~\text{rad} \;=\; 3.6976^\circ
  \label{eq:theta5_solution} \\
  \phi_{c,5} &= \pi - 5\theta_5 = 2.8189\ldots~\text{rad} \;=\; 161.51^\circ\,.
  \label{eq:phi_c5}
\end{align}

\subsection{Status: hypothesis, not theorem}

The solution~\eqref{eq:theta5_solution} is reported as a \emph{hypothesis}.
For $\mathbb{Z}_3$, the self-consistency condition was motivated by the
constancy of $Q_3$: ``the chirality breaking angle equals the unique
$\mathbb{Z}_3$ invariant.''
For $\mathbb{Z}_5$, no such unique invariant exists (Paper~VI).
The condition $5\theta_5 = \tilde{Q}_5(\theta_5)$ is a natural
generalisation, but a first-principles physical argument from the CDFT
vacuum equation of state would be needed to promote it to a theorem.


\section{Structural Comparison: $\mathbb{Z}_3$ vs $\mathbb{Z}_n$, $n \geq 5$}
\label{sec:comparison}

Table~\ref{tab:comparison} summarises the structural differences between
the lepton family and the higher families.

\begin{table}[h]
  \caption{Structural comparison: $\mathbb{Z}_3$ (leptons) vs $\mathbb{Z}_5$ (cinquefoil).}
  \label{tab:comparison}
  \begin{tabular}{lll}
    \toprule
    Property & $\mathbb{Z}_3$ (leptons) & $\mathbb{Z}_5$ (cinquefoil) \\
    \midrule
    Energy scale $M_n$ & $313.9$~MeV & $460.4$~MeV \\
    Number of modes & 3 & 5 \\
    Coefficient $c_n$ & $\sqrt{2}$ & $\sqrt{2}$ \\
    Valid domain (all $A_k > 0$) & Yes: $\theta \in (-\pi/12, \pi/12)$ & No \\
    $Q_n = 2/n$ constant? & Yes: $Q_3 = 2/3$ & No: $\tilde{Q}_5 \in [0.294, 0.328]$ \\
    $\sum m_k = 2nM_n$ & Yes: $1883.2$~MeV (verified) & Yes: $4603.9$~MeV (predicted) \\
    $\theta_n$ derivation & Algebraic (Paper~V) & Transcendental (hypothesis) \\
    Sector closed? & Yes (Paper~V) & Open \\
    Amplitude split & 3 positive always & $3+2$ or $4+1$ \\
    \bottomrule
  \end{tabular}
\end{table}

The key distinction: $\mathbb{Z}_3$ has a connected valid domain because
the minimum amplitude $\min_k A_k = 1 - \sqrt{2}\cos(0) + \ldots$
reaches a \emph{positive} maximum ($+0.293$ at $\theta = 0$).
For $n \geq 5$, the analogous maximum is \emph{negative}
(e.g.\ $-0.144$ for $n = 5$), so no valid domain exists.
This is the root of $\mathbb{Z}_3$ uniqueness (Paper~VI, Theorem~2).

The universal mass sum rule bridges the gap: it provides a firm prediction
for all $n$ without requiring the valid-domain condition.


\section{Open Problems}
\label{sec:open}

\begin{enumerate}

\item \textbf{Physical principle for $\theta_5$.}
  The transcendental equation~\eqref{eq:sc_5} gives a unique numerical
  $\theta_5 \approx 3.70^\circ$.
  A first-principles derivation from the CDFT vacuum equation of state ---
  analogous to how Paper~IV identified the functional form $\Phi_c(\rho_0)$
  as the remaining input --- is the priority open problem.

\item \textbf{Identify the cinquefoil mass spectrum.}
  Once $\theta_5$ is determined (or bounded), the five mode masses
  $m_k = M_5 A_k^2$ can be compared to known hadronic states
  at the $\sim 460$~MeV energy scale.
  The total mass $\Sigma_5 = 4603.9$~MeV is a testable sum rule.

\item \textbf{Physical meaning of the $3 + 2$ mode split.}
  The positive-amplitude modes ($A_k > 0$) and negative-amplitude modes
  ($A_k < 0$) interact differently with the vacuum background.
  Whether this split corresponds to a known quantum number
  (charge, colour, isospin) is an open question.

\item \textbf{Higher families $n \geq 7$.}
  The mass sum rule $\Sigma_n = 2nM_n$ applies to all $n$.
  The individual mode spectra and their identification with known or
  predicted particles require repeating the present analysis for each $n$.

\end{enumerate}


\section{Conclusion}

The universal mass sum rule $\sum_k m_k = 2nM_n$ is derived from
two algebraic facts: the $\mathbb{Z}_n$ Fourier identity
(which forces $\sum A_k^2 = 2n$ regardless of $\theta_n$) and the
Faddeev--Niemi energy scaling $M_n = M_3(n/3)^{3/4}$.
Together they give the $n^{7/4}$ mass hierarchy (Eq.~\eqref{eq:n74}),
verified for $n = 3$ and predicted for all higher families.

For $\mathbb{Z}_5$ specifically, we prove the $3 + 2$ mode split
(Theorem~\ref{thm:Z5_split}) and derive the transcendental
self-consistency equation~\eqref{eq:sc_5} for $\theta_5$,
reporting its solution $\theta_5 \approx 3.70^\circ$ as a hypothesis.

The mass sum $\Sigma_5 = 4603.9$~MeV is a model-level prediction, conditional
on the torus-knot hierarchy and the $m_e$ mass-unit calibration.


\section*{Supplementary Material}

This manuscript is accompanied by release-archive supplementary material in
the canonical Part I tree. The relevant paper-local Python scripts, numerical checks, generated figure outputs,
tabulated data, figure assets, and environment notes are maintained under
\texttt{notebooks/}, \texttt{outputs/}, \texttt{figures/}, and
\texttt{REPRODUCIBILITY.md} in \texttt{Part\_I\_Fundamental\_Physics/}.


\section*{Data and Code Availability}

The release-archive supplementary material is included in the archived Part I
release on Zenodo at \url{https://doi.org/10.5281/zenodo.20250820}. The current version DOI is assigned by Zenodo after each GitHub release. Zenodo is the
permanent release archive, and the public source archive is available on GitHub at
\url{https://github.com/bampita-bico/CDFD-Part-I-Release}. The broader public CDFD
Runtime is separate reusable software infrastructure; the numerical values reported
here are reproduced from this paper's Supplementary Material.


\section*{Acknowledgments}
The author used AI-assisted tools for formatting checks and code-verification
support. The scientific claims, interpretations, and final manuscript decisions
are the author's responsibility.

The author is indebted to the foundational work of
Y.~Koide~\cite{Koide1983} (lepton mass formula),
C.A.~Brannen~\cite{Brannen2006} (Koide parametrisation),
L.~Faddeev and A.J.~Niemi~\cite{FaddeevNiemi1997} (knots and particles),
and D.~Rolfsen~\cite{Rolfsen1976} (torus knot classification)
--- whose work underpins the mass sum rule derived here.

\section*{Declarations}

\subsection*{Funding}
This research received no external institutional funding.

\subsection*{Author Contributions}
Steve Bico Mujjabi: conceptualization, methodology, formal analysis, software,
validation, visualization, writing--original draft, and writing--review and
editing.

\subsection*{Conflicts of Interest}
The author declares no conflicts of interest.

\section*{References}
\bibliographystyle{plainnat}
\bibliography{references}

\end{document}
